% GENERATED by scripts/build-lecture-slides.py; do not edit. Portable profile: XeTeX, theme and fonts from the core.
\makeatletter\IfFileExists{core/beamer/beamerthememx.sty}{\def\mxcorepath{core/}\def\input@path{{core/beamer/}{core/}}}{\def\mxcorepath{../core/}\def\input@path{{../core/beamer/}{../core/}}}\makeatother
\documentclass[aspectratio=169,10pt,t]{beamer}
\usetheme{mx}
\usepackage{graphicx,booktabs}
\usepackage[french]{babel}
\setlength{\parskip}{.45em}
\hypersetup{pdfauthor={},pdfsubject={Lecture slides},pdfkeywords={},pdftitle={Généralités sur les fonctions numériques}}
\providecommand{\N}{\mathbb{N}}
\providecommand{\Z}{\mathbb{Z}}
\providecommand{\D}{\mathbb{D}}
\providecommand{\Q}{\mathbb{Q}}
\providecommand{\R}{\mathbb{R}}
\providecommand{\C}{\mathbb{C}}
\providecommand{\vect}[1]{\overrightarrow{#1}}
\providecommand{\vi}{\vec{\imath}}
\providecommand{\vj}{\vec{\jmath}}
\providecommand{\vk}{\vec{k}}
\providecommand{\norm}[1]{\left\lVert #1\right\rVert}
\providecommand{\vc}[2]{\begin{pmatrix}#1\\#2\end{pmatrix}}
\providecommand{\deter}[2]{\det\left(#1,#2\right)}
\providecommand{\abs}[1]{\left\lvert #1\right\rvert}
\providecommand{\dd}{\,\mathrm{d}}
\providecommand{\e}{\mathrm{e}}
\providecommand{\Card}{\operatorname{Card}}
\providecommand{\E}{\operatorname{E}}
\providecommand{\V}{\operatorname{V}}
\providecommand{\degre}{^{\circ}}
\providecommand{\un}[1]{\,\mathrm{#1}}
\providecommand{\sci}[2]{#1\times10^{#2}}
\providecommand{\pgcd}{\operatorname{PGCD}}
\providecommand{\ppcm}{\operatorname{PPCM}}
\title{Généralités sur les fonctions numériques}
\subtitle{Cours et exemples}
\date{}
\author{}
\begin{document}
\mxtitleframe
\begin{frame}{Objectifs}
Lire une image et un antécédent sur une courbe comme sur une expression,
  déterminer un ensemble de définition, établir qu'une fonction est paire,
  impaire ou périodique, et en tirer l'intervalle sur lequel il suffit de
  l'étudier.
\end{frame}
\begin{frame}{Fonction, image et antécédent}
Soit \(\mathcal{D}\) une partie de \(\mathbb{R}\). Définir une \textbf{fonction numérique} \(f\) sur
\(\mathcal{D}\), c'est associer à tout réel \(x\) de \(\mathcal{D}\) un unique réel, noté \(f(x)\).
On écrit
\[
  f : \begin{array}{ccc}
        \mathcal{D} & \longrightarrow & \mathbb{R} \\[0.15em]
        x   & \longmapsto     & f(x)
      \end{array}
\]
Le réel \(f(x)\) s'appelle l'\textbf{image} de \(x\) par \(f\). Réciproquement, si
\(y\) est un réel et si \(x\) est un élément de \(\mathcal{D}\) tel que \(f(x) = y\), on dit
que \(x\) est un \textbf{antécédent} de \(y\) par \(f\).
\end{frame}
\begin{frame}{Fonction, image et antécédent}
Un réel de \(\mathcal{D}\) a une image et une seule. En revanche, un réel \(y\) peut
avoir plusieurs antécédents, ou n'en avoir aucun : la fonction carré donne
deux antécédents à \(4\), et aucun à \(-1\).
\end{frame}
\begin{frame}{Fonction, image et antécédent}
Soit \(f\) la fonction définie sur \(\mathbb{R}\) par \(f(x) = x^{2} - 3x\).

Image de \(4\). On remplace : \(f(4) = 4^{2} - 3 \times 4 = 16 - 12 = 4\).
L'image de \(4\) est \(4\).

Antécédents de \(-2\). On résout l'équation \(f(x) = -2\), c'est-à-dire
\[
  x^{2} - 3x = -2
  \quad\Longleftrightarrow\quad x^{2} - 3x + 2 = 0 .
\]
Le discriminant vaut \(\Delta = 9 - 8 = 1 > 0\), d'où les deux solutions
\(x_{1} = \dfrac{3-1}{2} = 1\) et \(x_{2} = \dfrac{3+1}{2} = 2\). Le réel \(-2\)
admet donc exactement deux antécédents : \(1\) et \(2\).
\end{frame}
\begin{frame}{Fonction, image et antécédent (2)}
Antécédents de \(-3\). On résout \(x^{2} - 3x + 3 = 0\), dont le
discriminant vaut \(\Delta = 9 - 12 = -3 < 0\). Le réel \(-3\) n'a aucun
antécédent par \(f\).
\end{frame}
\begin{frame}{L'ensemble de définition}
Lorsqu'une fonction \(f\) est donnée par une expression, son \textbf{ensemble de
définition}, noté \(\mathcal{D}_{f}\), est l'ensemble des réels \(x\) pour lesquels cette
expression a un sens, c'est-à-dire pour lesquels le calcul de \(f(x)\) est
possible dans \(\mathbb{R}\).
\end{frame}
\begin{frame}{L'ensemble de définition}
Pour les expressions polynomiales, rationnelles et à racines carrées,
les contraintes suivantes déterminent le domaine :
\end{frame}
\begin{frame}{L'ensemble de définition (2)}
\begin{enumerate}
\setcounter{enumi}{0}
  \item \(x\) annule un \textbf{dénominateur} : le quotient \(\dfrac{A}{B}\)
        n'existe que si \(B \neq 0\) ;
  
\item \(x\) rend négatif ce qui se trouve \textbf{sous une racine carrée} :
        \(\sqrt{A}\) n'existe que si \(A \geq 0\) ;
  
\item les deux à la fois, lorsqu'une racine figure au dénominateur :
        \(\dfrac{1}{\sqrt{A}}\) n'existe que si \(A > 0\), l'inégalité devenant
        stricte.

\end{enumerate}
\end{frame}
\begin{frame}{L'ensemble de définition (3)}
Une fonction polynôme n'est jamais concernée : son ensemble de définition
est \(\mathbb{R}\) tout entier.
\end{frame}
\begin{frame}{L'ensemble de définition}
\begin{enumerate}
\setcounter{enumi}{0}\item Repérer dans l'expression tous les dénominateurs et toutes les
        racines carrées. Pour un polynôme, conclure : \(\mathcal{D}_f=\mathbb{R}\). Pour une fonction trigonométrique, vérifier aussi son domaine; la tangente exige \(\cos x\ne0\).
  
\item Écrire une condition par obstacle rencontré : une équation
        \(B \neq 0\), une inéquation \(A \geq 0\).

\end{enumerate}
\end{frame}
\begin{frame}{L'ensemble de définition (2)}
\begin{enumerate}
\setcounter{enumi}{2}\item Résoudre chaque condition séparément. Un tableau de signes est
        souvent le chemin le plus court pour une inéquation du second degré
        ou pour un quotient.
  
\item Conserver les réels qui vérifient \emph{toutes} les conditions à la
        fois : l'ensemble de définition est leur intersection.
  
\item Écrire la réponse sous forme d'intervalles ou de réunion
        d'intervalles, jamais sous forme de phrase.

\end{enumerate}
\end{frame}
\begin{frame}{L'ensemble de définition}
Soit \(f(x) = \dfrac{x+5}{x^{2} - x - 6}\).

Le seul obstacle est le dénominateur : il faut \(x^{2} - x - 6 \neq 0\). On
résout d'abord l'équation associée : \(\Delta = 1 + 24 = 25\), donc
\(x_{1} = \dfrac{1-5}{2} = -2\) et \(x_{2} = \dfrac{1+5}{2} = 3\).

Ces deux réels, et eux seuls, sont à exclure :
\[
  \mathcal{D}_{f} = \mathbb{R} \setminus \{-2\,;\,3\}
          = \left]-\infty\,;\,-2\right[ \cup \left]-2\,;\,3\right[
            \cup \left]3\,;\,+\infty\right[.
\]
\end{frame}
\begin{frame}{L'ensemble de définition}
Soit \(g(x) = \dfrac{\sqrt{x+1}}{x-4}\).

Deux obstacles cette fois. Sous la racine : \(x + 1 \geq 0\), soit
\(x \geq -1\). Au dénominateur : \(x - 4 \neq 0\), soit \(x \neq 4\).

On garde les réels vérifiant les deux conditions :
\[
  \mathcal{D}_{g} = \left[-1\,;\,4\right[ \cup \left]4\,;\,+\infty\right[.
\]
Remarquez que \(-1\) est conservé - la racine y vaut \(0\), ce qui est
parfaitement licite - tandis que \(4\) est exclu.
\end{frame}
\begin{frame}{Sens de variation}
Soit \(f\) une fonction définie sur un intervalle \(I\).

\begin{itemize}

  \item \(f\) est \textbf{croissante} sur \(I\) lorsque, pour tous réels \(a\) et
        \(b\) de \(I\) : \(a \leq b \implies f(a) \leq f(b)\).
  
\item \(f\) est \textbf{décroissante} sur \(I\) lorsque, pour tous réels \(a\) et
        \(b\) de \(I\) : \(a \leq b \implies f(a) \geq f(b)\).

\end{itemize}
\end{frame}
\begin{frame}{Sens de variation (2)}
Pour la croissance stricte, \(a<b\) implique \(f(a)<f(b)\); pour la
décroissance stricte, \(a<b\) implique \(f(a)>f(b)\). Une fonction \textbf{monotone} sur \(I\) est une fonction
qui y est croissante, ou qui y est décroissante.
\end{frame}
\begin{frame}{Sens de variation}
La fonction inverse est décroissante sur \(\left]-\infty\,;\,0\right[\) et
décroissante sur \(\left]0\,;\,+\infty\right[\), mais elle n'est pas décroissante
sur leur réunion. Prenons en effet \(a = -1\) et \(b = 1\) : on a bien
\(a \leq b\), et pourtant \(f(a) = -1\) tandis que \(f(b) = 1\), de sorte que
\(f(a) < f(b)\) - l'inégalité n'a pas été renversée. Dire « \(f\) est
décroissante sur \(\mathbb{R}^{*}\) » est donc faux. On écrit deux lignes, pas une.
\end{frame}
\begin{frame}{Sens de variation}
Soit \(f\) une fonction monotone sur un intervalle \(I\) et \(\lambda\) un réel non
nul. La fonction \(\lambda f\) a le même sens de variation que \(f\) sur \(I\) si
\(\lambda > 0\), et le sens contraire si \(\lambda < 0\).
\end{frame}
\begin{frame}{Sens de variation}
Supposons \(f\) croissante sur \(I\) et \(\lambda > 0\). Soient \(a\) et \(b\) deux
réels de \(I\) tels que \(a \leq b\). Par croissance de \(f\), on a
\(f(a) \leq f(b)\). Multiplier les deux membres d'une inégalité par un réel
strictement positif en conserve le sens, donc
\(\lambda f(a) \leq \lambda f(b)\), c'est-à-dire
\((\lambda f)(a) \leq (\lambda f)(b)\) : la fonction \(\lambda f\) est croissante
sur \(I\).
\end{frame}
\begin{frame}{Sens de variation (2)}
Si \(\lambda < 0\), la même multiplication renverse l'inégalité et donne
\(\lambda f(a) \geq \lambda f(b)\) : la fonction \(\lambda f\) est alors
décroissante sur \(I\). Le cas où \(f\) est décroissante se traite mot pour mot
de la même manière.
\end{frame}
\begin{frame}{Parité : quand la courbe se reflète}
Une partie \(\mathcal{D}\) de \(\mathbb{R}\) est \textbf{symétrique par rapport à zéro} lorsque,
pour tout réel \(x\), l'appartenance de \(x\) à \(\mathcal{D}\) entraîne celle de \(-x\) à
\(\mathcal{D}\).
\end{frame}
\begin{frame}{Parité : quand la courbe se reflète}
Soit \(f\) une fonction dont l'ensemble de définition \(\mathcal{D}_{f}\) est symétrique
par rapport à zéro.

\begin{itemize}

  \item \(f\) est \textbf{paire} lorsque \(f(-x) = f(x)\) pour tout \(x\) de
        \(\mathcal{D}_{f}\).
  
\item \(f\) est \textbf{impaire} lorsque \(f(-x) = -f(x)\) pour tout \(x\) de
        \(\mathcal{D}_{f}\).

\end{itemize}
\end{frame}
\begin{frame}{Parité : quand la courbe se reflète}
Le plan est rapporté à un repère orthogonal \(\left(O\,;\,\overrightarrow{\imath},\,\overrightarrow{\jmath}\right)\) et \(\mathcal{C}_{f}\) désigne la
courbe représentative de \(f\).

\begin{itemize}

  \item \(f\) est paire si et seulement si \(\mathcal{C}_{f}\) est symétrique par rapport à l'axe des
        ordonnées.
  
\item \(f\) est impaire si et seulement si \(\mathcal{C}_{f}\) est symétrique par rapport à
        l'origine \(O\) du repère.

\end{itemize}
\end{frame}
\begin{frame}{Parité : quand la courbe se reflète}
Le point \(M\) de coordonnées \((x\,;\,y)\) a pour symétrique par rapport à l'axe
des ordonnées le point \(M'\) de coordonnées \((-x\,;\,y)\).
\end{frame}
\begin{frame}{Parité : quand la courbe se reflète (2)}
Supposons \(f\) paire et prenons un point \(M(x\,;\,f(x))\) de \(\mathcal{C}_{f}\). Son
symétrique est \(M'(-x\,;\,f(x))\). Comme \(\mathcal{D}_{f}\) est symétrique par rapport
à zéro, le réel \(-x\) appartient à \(\mathcal{D}_{f}\), et la parité donne
\(f(-x) = f(x)\). L'ordonnée de \(M'\) est donc bien \(f(-x)\) : le point \(M'\)
appartient à \(\mathcal{C}_{f}\). La courbe est donc globalement invariante par cette
symétrie.
\end{frame}
\begin{frame}{Parité : quand la courbe se reflète (3)}
Réciproquement, si \(\mathcal{C}_{f}\) est symétrique par rapport à l'axe des
ordonnées, alors pour tout \(x\) de \(\mathcal{D}_{f}\) le symétrique
\(M'(-x\,;\,f(x))\) du point \(M(x\,;\,f(x))\) appartient à \(\mathcal{C}_{f}\) ; son
ordonnée vaut donc \(f(-x)\), d'où \(f(-x) = f(x)\).
\end{frame}
\begin{frame}{Parité : quand la courbe se reflète}
\begin{enumerate}
\setcounter{enumi}{0}
  \item Déterminer \(\mathcal{D}_{f}\) et vérifier qu'il est symétrique par rapport à
        zéro. Si ce n'est pas le cas, conclure immédiatement : la fonction
        n'est ni paire ni impaire, et il n'y a rien d'autre à écrire.
  
\item Calculer \(f(-x)\) en remplaçant partout \(x\) par \(-x\), puis simplifier
        \emph{complètement} l'expression obtenue.
  
\item Comparer le résultat à \(f(x)\) et à \(-f(x)\) :

\end{enumerate}
\end{frame}
\begin{frame}{Parité : quand la courbe se reflète (2)}
\begin{itemize}

          \item si \(f(-x) = f(x)\), la fonction est paire ;
          
\item si \(f(-x) = -f(x)\), la fonction est impaire ;
          
\item si aucune des deux égalités n'est vraie pour tout \(x\), elle
                n'est ni l'une ni l'autre - et un seul contre-exemple
                numérique pour chacune des deux égalités suffit alors à le prouver.

\end{itemize}
\end{frame}
\begin{frame}{Parité : quand la courbe se reflète}
\textbf{a)} \(f(x) = 3x^{4} - x^{2} + 5\), définie sur \(\mathbb{R}\), qui est symétrique
par rapport à zéro. On calcule
\[
  f(-x) = 3(-x)^{4} - (-x)^{2} + 5 = 3x^{4} - x^{2} + 5 = f(x).
\]
La fonction \(f\) est paire. C'était prévisible : tous les exposants de \(x\)
sont pairs, et la constante \(5\) se lit comme \(5x^{0}\).

\textbf{b)} \(g(x) = \dfrac{x}{x^{2}+1}\), définie sur \(\mathbb{R}\). On calcule
\[
  g(-x) = \dfrac{-x}{(-x)^{2}+1} = \dfrac{-x}{x^{2}+1}
        = -\,\dfrac{x}{x^{2}+1} = -g(x).
\]
La fonction \(g\) est impaire.
\end{frame}
\begin{frame}{Parité : quand la courbe se reflète (2)}
\textbf{c)} \(h(x) = x^{2} + x\), définie sur \(\mathbb{R}\). On calcule
\(h(-x) = x^{2} - x\). Pour conclure, un contre-exemple suffit : \(h(1) = 2\) et
\(h(-1) = 0\). Comme \(h(-1) \neq h(1)\) la fonction n'est pas paire, et comme
\(h(-1) \neq -h(1)\) elle n'est pas impaire non plus.
\end{frame}
\begin{frame}{Parité : quand la courbe se reflète}
Soit \(f\) une fonction et \(\mathcal{C}_{f}\) sa courbe dans un repère orthogonal
\(\left(O\,;\,\overrightarrow{\imath},\,\overrightarrow{\jmath}\right)\). Soient \(a\) et \(b\) deux réels. Le domaine doit être symétrique par rapport à \(a\) : \(a+h\) appartient au domaine si et seulement si \(a-h\) y appartient. Sous cette condition,
\end{frame}
\begin{frame}{Parité : quand la courbe se reflète (2)}
\begin{itemize}

  \item La droite d'équation \(x = a\) est un \textbf{axe de symétrie} de
        \(\mathcal{C}_{f}\) si et seulement si, pour tout réel \(h\) tel que \(a+h\) et \(a-h\)
        appartiennent à \(\mathcal{D}_{f}\),
        \[ f(a+h) = f(a-h). \]
  
\item Le point \(\Omega(a\,;\,b)\) est un \textbf{centre de symétrie} de
        \(\mathcal{C}_{f}\) si et seulement si, pour tout réel \(h\) tel que \(a+h\) et \(a-h\)
        appartiennent à \(\mathcal{D}_{f}\),
        \[ f(a+h) + f(a-h) = 2b. \]

\end{itemize}
\end{frame}
\begin{frame}{Parité : quand la courbe se reflète}
Le symétrique du point \(M(a+h\,;\,f(a+h))\) par rapport au point
\(\Omega(a\,;\,b)\) est le point \(M'\) tel que \(\Omega\) soit le milieu de
\([MM']\). En écrivant les coordonnées du milieu, \(M'\) a pour abscisse
\(2a - (a+h) = a - h\) et pour ordonnée \(2b - f(a+h)\).

Dire que \(\Omega\) est centre de symétrie de \(\mathcal{C}_{f}\), c'est dire que \(M'\)
appartient à \(\mathcal{C}_{f}\), c'est-à-dire que son ordonnée est l'image de son
abscisse :
\[
  2b - f(a+h) = f(a-h),
  \quad\text{soit}\quad
  f(a+h) + f(a-h) = 2b .
\]
\end{frame}
\begin{frame}{Parité : quand la courbe se reflète}
En prenant \(a = 0\), la première égalité devient \(f(h) = f(-h)\) : c'est la
parité. En prenant \(a = 0\) et \(b = 0\), la seconde devient
\(f(h) + f(-h) = 0\) : c'est l'imparité. Il n'y a donc qu'un seul résultat à
retenir, la parité n'en étant que la version centrée à l'origine.
\end{frame}
\begin{frame}{Périodicité : quand la courbe se répète}
Soit \(f\) une fonction définie sur \(\mathcal{D}_{f}\) et \(T\) un réel strictement
positif. La fonction \(f\) est périodique de période \(T\) lorsque, pour
tout \(x\) de \(\mathcal{D}_{f}\), les réels \(x+T\) et \(x-T\) appartiennent à \(\mathcal{D}_{f}\) et
\[
  f(x+T) = f(x).
\]
Le plus petit réel strictement positif vérifiant cette égalité, lorsqu'il
existe, s'appelle la \textbf{période} de \(f\).
\end{frame}
\begin{frame}{Périodicité : quand la courbe se répète}
Si \(f\) est périodique de période \(T\), alors pour tout entier relatif \(k\) et
tout \(x\) de \(\mathcal{D}_{f}\) :
\[
  f(x + kT) = f(x).
\]
Il suffit donc d'étudier \(f\) sur un intervalle de longueur \(T\) - par exemple
\(\left[0\,;\,T\right[\) - et de reproduire la portion de courbe obtenue par
translations successives de vecteur \(T\overrightarrow{\imath}\) pour obtenir la courbe entière.
\end{frame}
\begin{frame}{Périodicité : quand la courbe se répète}
L’égalité est vraie pour \(k=0\). Si elle est vraie pour un entier \(k\ge0\), alors
\(f(x+(k+1)T)=f((x+kT)+T)=f(x+kT)=f(x)\); une récurrence donne tous les entiers naturels.
En appliquant la définition en \(x-T\), on obtient \(f(x-T)=f(x)\). Répéter ce déplacement donne les entiers négatifs. Le domaine reste invariant à chaque déplacement.
\end{frame}
\begin{frame}{Périodicité : quand la courbe se répète}
Les fonctions \(x \mapsto \cos x\) et \(x \mapsto \sin x\) sont définies sur \(\mathbb{R}\)
et périodiques de période \(2\pi\) : un tour complet du cercle
trigonométrique ramène au même point, donc aux mêmes coordonnées.

La fonction \(f(x) = \cos(3x)\) est également périodique. Cherchons sa période.
Pour tout réel \(x\),
\[
  f\!\left(x + \dfrac{2\pi}{3}\right)
  = \cos\!\left(3x + 2\pi\right) = \cos(3x) = f(x),
\]
donc \(\dfrac{2\pi}{3}\) est une période de \(f\).
\end{frame}
\begin{frame}{Périodicité : quand la courbe se répète}
Lorsqu'une fonction est à la fois périodique et paire (ou impaire), les deux
propriétés se cumulent et l'intervalle d'étude se réduit deux fois.
\end{frame}
\begin{frame}{Périodicité : quand la courbe se répète (2)}
\begin{enumerate}
\setcounter{enumi}{0}\item La périodicité de période \(T\) ramène l'étude à un intervalle de
        longueur \(T\), centré si possible en zéro :
        \(\left[-\dfrac{T}{2}\,;\,\dfrac{T}{2}\right]\).
  
\item La parité, ou l'imparité, permet alors de n'étudier que la moitié
        droite : \(\left[0\,;\,\dfrac{T}{2}\right]\).

\end{enumerate}
\end{frame}
\begin{frame}{Périodicité : quand la courbe se répète (3)}
\begin{enumerate}
\setcounter{enumi}{2}\item On trace la courbe sur ce dernier intervalle, on complète par la
        symétrie donnée par la parité, puis on reproduit le motif obtenu par
        translations de vecteur \(T\overrightarrow{\imath}\).

\end{enumerate}
\end{frame}
\begin{frame}{Périodicité : quand la courbe se répète}
Soit \(f(x) = \cos(2x)\), définie sur \(\mathbb{R}\).

\textbf{Périodicité.} Pour tout réel \(x\), \(f(x + \pi) = \cos(2x + 2\pi)
= \cos(2x) = f(x)\). La fonction est périodique de période \(\pi\), et l'on
peut se limiter à \(\left[-\dfrac{\pi}{2}\,;\,\dfrac{\pi}{2}\right]\).

\textbf{Parité.} Pour tout réel \(x\), \(f(-x) = \cos(-2x) = \cos(2x) = f(x)\),
car la fonction cosinus est paire. La fonction \(f\) est donc paire, et l'on
peut se limiter à \(\left[0\,;\,\dfrac{\pi}{2}\right]\).
\end{frame}
\begin{frame}{Périodicité : quand la courbe se répète (2)}
D'un intervalle infini, on est passé à un intervalle de longueur
\(\dfrac{\pi}{2}\).
\end{frame}
\end{document}
